% Manuscript: exact label-setting DP for the four-objective lot-sizing model of Kania et al.
% Target: European Journal of Operational Research (subscription route, no publication fee).
% Build (from the paper/ folder, after `python3 make_figures.py --suite ../suite --out .`):
%     pdflatex main && bibtex main && pdflatex main && pdflatex main
\documentclass[preprint,12pt,authoryear]{elsarticle}
\usepackage[T1]{fontenc}
\usepackage{lmodern}
\usepackage{amsmath,amssymb,amsthm,mathtools}
\usepackage{booktabs}
\usepackage{graphicx}
\usepackage{enumitem}
\usepackage[hidelinks]{hyperref}
\usepackage[capitalise,noabbrev]{cleveref}
% Theorem environments and notation shared by main.tex and the appendix.
\theoremstyle{plain}
\newtheorem{theorem}{Theorem}[section]
\newtheorem{lemma}[theorem]{Lemma}
\newtheorem{proposition}[theorem]{Proposition}
\newtheorem{corollary}[theorem]{Corollary}
\theoremstyle{definition}
\newtheorem{assumption}{Assumption}
\renewcommand{\theassumption}{A\arabic{assumption}}
\newtheorem{definition}[theorem]{Definition}
\theoremstyle{remark}
\newtheorem{remark}[theorem]{Remark}

\newcommand{\Tp}{T'}
\newcommand{\cQ}{\mathcal{Q}}
\newcommand{\cL}{\mathcal{L}}
\newcommand{\cumD}{\mathrm{cD}}
\newcommand{\Ibar}{\bar I}
\newcommand{\POC}{\mathrm{POC}}
\newcommand{\HC}{\mathrm{HC}}
\newcommand{\CSL}{\mathrm{CSL}}
\newcommand{\ITO}{\mathrm{ITO}}
\newcommand{\Ss}{\mathrm{SS}}
\newcommand{\SOT}{\mathrm{SOT}}
\newcommand{\MS}{\mathrm{MS}}
\newcommand{\moq}{\mathrm{moq}}
\newcommand{\eps}{\varepsilon}
\crefname{assumption}{Assumption}{Assumptions}
\newcommand{\lceilL}[1]{\lceil #1\rceil_{\cL}}


\IfFileExists{tables/software.tex}{\newcommand{\verCpu}{12th Gen Intel(R) Core(TM) i5-1235U}
\newcommand{\verOs}{Ubuntu 26.04.1 LTS}
\newcommand{\verGcc}{g++ (Ubuntu 15.2.0-16ubuntu1) 15.2.0}
\newcommand{\verCmake}{4.2.3}
\newcommand{\verNinja}{1.13.2}
\newcommand{\verPython}{3.14.4}
\newcommand{\verNumpy}{2.5.5}
\newcommand{\verMatplotlib}{3.11.2}
\newcommand{\verScipy}{1.18.1}
\newcommand{\verPymoo}{0.6.2}
\newcommand{\verHighs}{HiGHS (bundled with SciPy)}
}{}
\IfFileExists{tables/e3_full.tex}{\newcommand{\EfullTime}{8.5}
\newcommand{\EfullFront}{792\,232}
}{\newcommand{\EfullTime}{??}\newcommand{\EfullFront}{??}}
\IfFileExists{tables/e6_macros.tex}{\newcommand{\EsixInst}{700}
\newcommand{\EsixSolved}{700}
\newcommand{\EsixSeries}{300}
\newcommand{\EsixChkOK}{30}\newcommand{\EsixChkN}{30}
\newcommand{\EsixMedTime}{0.20}
\newcommand{\EsixMaxTime}{59.7}
\newcommand{\EsixTotalMin}{10}
\newcommand{\EsixTimeout}{0}
\newcommand{\EsixTlim}{30}
\newcommand{\EsixPctSolved}{100.0}
}{\newcommand{\EsixInst}{??}\newcommand{\EsixSolved}{??}\newcommand{\EsixSeries}{??}\newcommand{\EsixChkOK}{??}\newcommand{\EsixChkN}{??}\newcommand{\EsixMedTime}{??}\newcommand{\EsixMaxTime}{??}\newcommand{\EsixTotalMin}{??}\newcommand{\EsixTimeout}{??}\newcommand{\EsixTlim}{??}\newcommand{\EsixPctSolved}{??}}
\IfFileExists{tables/e4_inv.tex}{\newcommand{\EinvOK}{5}\newcommand{\EinvN}{5}
}{\newcommand{\EinvOK}{??}\newcommand{\EinvN}{??}}
\newcommand{\tabin}[1]{\IfFileExists{tables/#1}{\input{tables/#1}}{\fbox{run make\_figures.py: #1}}}
\newcommand{\figin}[2]{\IfFileExists{figures/#2}{\includegraphics[width=#1]{figures/#2}}{\fbox{run make\_figures.py: #2}}}

\journal{European Journal of Operational Research}

\begin{document}
\begin{frontmatter}

\title{Exact Pareto fronts for a four-objective lot-sizing model with safety stock and
safety order time: a label-setting dynamic programming approach}

\author[unimed]{Michael Sipayung\corref{cor}}
\ead{EMAIL@unimed.ac.id}
\cortext[cor]{Corresponding author.}
\affiliation[unimed]{organization={Department of Physics, Faculty of Mathematics and Natural
Sciences, Universitas Negeri Medan},
  city={Medan}, country={Indonesia}}

\begin{abstract}
Interactive multiobjective methods for lot sizing let a decision maker navigate among
Pareto optimal production plans, but in the published models the plans are produced by
approximate solvers, so the decision maker never knows how far the shown trade-offs are
from the attainable ones. We study the single-item, multi-period model of Kania et al.\
(2022), which integrates lot sizing with the placement of a safety stock and a safety
order time and optimises purchasing and ordering cost, holding cost, cycle service level
and inventory turnover simultaneously. We show that, for fixed safety parameters, the
model is a multi-criteria shortest path problem on a layered graph and develop an exact
label-setting dynamic program. We prove that every Pareto-optimal order raises the
cumulative quantity to the level that covers demand until the next order, an analogue of
the Wagner--Whitin property, which reduces the branching from order quantities to order
periods. Together with an exact memoization over classes of rounded requirement vectors
and a dominance rule between classes, this makes the algorithm robust to demand that
exceeds its forecast. For continuous safety parameters we prove
that the vertices of a line arrangement suffice to obtain the complete front. On two
industrial items with horizons of 48 weeks, the algorithm computes complete fronts of
$2\times10^4$ to $7\times10^4$ points in seconds on a laptop. An $\varepsilon$-constraint
mixed-integer programme solved with HiGHS reproduces the same fronts more than four orders
of magnitude more slowly, and NSGA-III recovers less than one per cent of the exact front
even with a problem-specific repair operator. On a public benchmark of \EsixInst{}
instances built from the M5 (Walmart) and UCI Online Retail II demand data, covering
smooth, erratic, intermittent and lumpy demand, \EsixPctSolved\% of the instances are solved
exactly within \EsixTlim{} minutes, with a median time of \EsixMedTime{} seconds, and the
remaining ones are explained by the complexity analysis. The exact fronts make interactive lot sizing
verifiable and provide reference sets for learning-based solvers.
\end{abstract}

\begin{keyword}
Multiple objective programming \sep Lot sizing \sep Dynamic programming \sep
Multi-criteria shortest path \sep Safety stock
\end{keyword}
\end{frontmatter}

%=====================================================================================
\section{Introduction}\label{sec:intro}

Lot sizing decides how much to order and when, and it has been a core problem of
production planning since the dynamic programme of \citet{WagnerWhitin1958}. Most of the
literature optimises a single cost function \citep{Brahimi2017}, but practitioners
balance several goals that cannot be expressed in one currency: purchasing and ordering
costs, the capital tied up in stock, the probability of serving demand from stock and the
speed at which inventory turns over. When demand is uncertain these goals conflict
sharply, because the safety stock that protects service is exactly the stock that raises
holding cost and lowers turnover.

Kania and co-authors addressed this conflict with a sequence of multiobjective models and
interactive solution processes \citep{Kania2021ICCL,Kania2022CIE,Kania2024DESMILS,
Kania2023thesis}. Their model of \citet{Kania2022CIE}, which is the subject of this paper,
plans the order quantities of one item over a horizon of weekly periods and places two
safety instruments at the same time: a safety stock and a \emph{safety order time}, which
expresses the additional stock as a multiple of the period demand. Four objectives are
optimised: purchasing and ordering cost, holding cost, cycle service level and inventory
turnover. The decision maker explores the trade-offs with the trade-off-free interactive
method E-NAUTILUS \citep{Ruiz2015ENautilus}, which navigates over a precomputed set of
Pareto optimal solutions. Related work by \citet{Heikkinen2023} follows the same
interactive philosophy with forecast-driven demand.

The quality of such a process depends on the quality of the solutions it navigates. In
the published applications the set is approximated by a general-purpose solver, and
nothing guarantees that the plans offered to the decision maker are Pareto optimal or
that the set covers the attainable trade-offs. A decision maker who settles for a plan
that is in fact dominated pays for the approximation without knowing it. Whether the
\emph{exact} front of this model can be computed fast enough to support interactive use
has not been studied.

This paper answers the question affirmatively. Our first contribution is structural:
once the safety parameters are fixed, a cumulative reformulation turns the inner problem
into a multi-criteria shortest path problem on a layered acyclic graph whose arc costs
are additive in the number of orders, the cumulative quantity and the turnover terms. The
second contribution is an exact label-setting dynamic programme together with three
reductions whose safety we prove: an upper bound on the total order quantity, a rule that
removes orders placed before they are needed, and a memoization that solves one inner
problem per class of rounded requirement vectors and keeps only the largest service
level per class. The third contribution handles continuous safety parameters, for which
we prove that the vertices of an arrangement of lines carry the complete Pareto front;
this replaces an infinite outer domain by a finite candidate set without loss. The fourth
contribution is a computational study on the two industrial items of \citet{Kania2022CIE}
that verifies exactness against exhaustive enumeration and an independent mixed-integer
formulation, measures the empirical complexity, and quantifies how much of the front a
state-of-the-art evolutionary algorithm actually recovers. The fifth contribution is a
public benchmark of lot-sizing instances derived from real demand data, with fixed and
documented rules for every parameter, together with their exact Pareto fronts as reference
sets for future exact and approximate methods. All code, instances and scripts are public.

The remainder is organised as follows. \Cref{sec:related} reviews related work.
\Cref{sec:model} recalls the model, and \cref{sec:method} presents the algorithm, whose
proofs are collected in \cref{app:theory}. \Cref{sec:design} describes the experimental
design and \cref{sec:results} the results. \Cref{sec:discussion} discusses implications
and limitations, and \cref{sec:conclusion} concludes.

%=====================================================================================
\section{Related work}\label{sec:related}

Single-item dynamic lot sizing with a single cost objective is polynomially solvable by
dynamic programming in its classical form, and a large body of work extends it to
capacities, backlogging, remanufacturing and stochastic demand \citep{Brahimi2017}.
Multiobjective variants are fewer. The models of Kania and co-authors add service and
turnover objectives and safety instruments to the classical cost objective
\citep{Kania2021ICCL,Kania2022CIE}, and DESMILS extends the approach to many items by
solving the single-item problem for a few representative items only
\citep{Kania2024DESMILS}. These works concentrate on modelling and on the interactive
process. The solution sets are produced by population-based or scalarisation-based
approximation, and their distance to the true front is not reported.

Exact methods for multiobjective combinatorial problems are surveyed by
\citet{Ehrgott2005}. Two families are relevant here. Scalarisation methods such as the
$\varepsilon$-constraint method \citep{Haimes1971} and its augmented variants
\citep{Mavrotas2009} solve a sequence of single-objective mixed-integer programmes. They
are generic but require one solve per front point and become slow when the front is
large. Labelling algorithms for the multiobjective shortest path problem, from the early
work of \citet{Hansen1980} and \citet{Martins1984} to the biobjective and multiobjective
Dijkstra algorithms of \citet{SedenoNoda2019} and \citet{Maristany2021} and the
label-correcting method of \citet{Kergosien2022}, compute complete sets of efficient
paths by propagating non-dominated labels. They are efficient when the problem has a
graph structure with additive costs, which a lot-sizing model does not obviously have:
the turnover objective is a sum of reciprocals of average stock, and the service-level
objective depends on the safety parameters rather than on the plan.

Metaheuristics, in particular NSGA-III \citep{DebJain2014}, are the default tool for
many-objective problems in applied work, and implementations such as pymoo
\citep{BlankDeb2020} make them easy to use. Their weakness is the absence of any
guarantee; without an exact reference, a front that is far from optimal can look
convincing \citep{Sorensen2015}. Exact fronts are also the natural training and test
references for learning-based solvers such as Pareto set learning \citep{Lin2022PSL}. We
return to this point in \cref{sec:discussion}.

%=====================================================================================
\section{The model}\label{sec:model}

We use the model of \citet{Kania2022CIE} in the notation of \cref{app:model}, to which
we refer for all definitions. An item is planned over periods $t=1,\dots,T$ with
deterministic lead time $L$, so that orders are placed in the decision periods
$\tau=1,\dots,\Tp$ with $\Tp=T-L$. Order quantities belong to
$\cQ=\{0\}\cup\{\moq+a r:a\in\mathbb{N}_0\}$ with a minimum order quantity $\moq$ and a
rounding value $r$. The safety stock $\Ss$ and the safety order time $\SOT$ lie in the
domain $\Omega$ defined by $\SOT+\Ss/\mu\le\MS$, where $\mu$ is the mean period demand
and $\MS$ a managerial bound. The stock at the end of every period after the lead time
must cover $\Ss+\SOT\cdot D(t)$. The objectives are the purchasing and ordering cost
$\POC$ and the holding cost $\HC$ (both minimised), and the cycle service level
$\CSL=\Phi\bigl((\Ss+\mu\SOT)/\sigma\bigr)$ and the inventory turnover $\ITO$, defined as
the sum over periods of $(D(t)+\sigma)$ divided by the average stock (both maximised).
Turnover is defined only for positive average stock, so we require an average stock of
at least $\varepsilon$ units in every period. The domain of the safety parameters is
either a finite grid, as in the original study, or the full continuous set $\Omega$.

%=====================================================================================
\section{Exact algorithm}\label{sec:method}

\paragraph{Cumulative reformulation}
Writing the plan in terms of cumulative quantities $X(\tau)=\sum_{s\le\tau}Q(s)$ turns
every constraint into a lower bound $X(\tau)\ge\beta_\tau$ that is affine in the safety
parameters, and every objective into a function of consecutive pairs
$(X(\tau-1),X(\tau))$ (\cref{app:model}). For fixed safety parameters the service level
is constant, and a plan is a monotone path in a layered acyclic graph whose nodes are
pairs of a period and a cumulative quantity. Each arc carries three additive costs: one
if an order is placed, the cumulative quantity, and the turnover term of the period. The
three remaining objectives are strictly monotone functions of these sums at the
terminal node. The inner problem is therefore a multi-criteria shortest path problem
(\cref{app:inner}).

\paragraph{Label-setting dynamic programme}
We process the layers in order and keep, at each node, the set of non-dominated labels,
that is, triples of order count, cumulative-sum total and accumulated turnover. Because
arc existence and costs depend only on the arc, dominance at a node is preserved by every
extension, and the procedure returns the exact front of the inner problem
(\cref{thm:exact}). Labels are stored as one two-dimensional staircase per order count,
which makes the dominance test logarithmic.

\paragraph{Order-event programme}
The period-by-period graph lets an order raise the cumulative quantity to any reachable
value, so a node has up to $S$ successors. We prove that this freedom is never needed: in a
Pareto-optimal plan, an order raises the cumulative quantity to the smallest level that
satisfies the requirements and positivity conditions of all periods until the next order,
or to the minimum order quantity if that is larger (\cref{lem:cover}). The only exceptions
are a few levels at which moving one case pack to the next order would empty the stock in
that period. The proof is an exchange argument in the spirit of the Wagner--Whitin
property \citep{WagnerWhitin1958}: excess units can be moved to the next order without
changing the number of orders, which strictly reduces holding cost and does not reduce
turnover. The algorithm therefore works on a graph of order events, in which each node
branches over at most $\Tp$ periods of the next order, and remains exact
(\cref{thm:event}).

\paragraph{Provably safe reductions}
Three reductions keep the graph small. First, every plan whose total quantity exceeds the
largest requirement by more than one order step plus $2\varepsilon$ is dominated
(\cref{lem:ub}); this bounds the state space. Second, an order placed in a period where
it is not yet needed can be postponed to the next period with a strict improvement in
holding cost and no loss in the other objectives, provided the postponed plan keeps
positive average stock in both periods it affects (\cref{lem:D}). Third, the inner
problem depends on the safety parameters only through the vector of requirements rounded
up to reachable quantities (\cref{lem:memo}). Hence one inner problem is solved per
distinct vector, and only the largest service level of each class is kept
(\cref{cor:memo}). Moreover, a class whose rounded requirements are componentwise at least
those of another class with a service level at least as high cannot contribute to the
front and is skipped (\cref{cor:classdom}). The remaining classes are independent and are
solved in parallel.

\paragraph{Continuous safety parameters}
Each rounded requirement changes only when an affine function of $(\Ss,\SOT)$ crosses a
reachable quantity, so the continuous domain is partitioned by an arrangement of lines.
We prove that every point of $\Omega$ is weakly dominated by a vertex of this arrangement
whose rounded requirement vector is componentwise smaller and whose service level is at
least as high (\cref{thm:vertex}). The complete front over the continuous domain is
therefore obtained from finitely many candidate points. In our instances these number a
few hundred to a few thousand.

\paragraph{Ordering-cost invariance}
Because cumulative quantities are multiples of the rounding value, a sufficiently small
ordering cost cannot reverse the comparison between two plans with different total
quantities. The Pareto-optimal decisions are then independent of the ordering cost
(\cref{prop:c}), which explains a robustness that we also observe numerically.

\paragraph{Complexity}
With $S=\Theta(\mu T/r)$ states per layer and at most $O(T^2S)$ labels per node, one inner
problem of the period-by-period programme takes $O\bigl(T^{6}(\mu/r)^{3}\log(T\mu/r)\bigr)$
time in the worst case, and the order-event programme improves this by a factor of order
$S/T$ (\cref{app:complexity}). The bound is pseudo-polynomial and, as \cref{sec:results} shows,
predicts the observed growth closely.

\paragraph{Implementation}
The algorithm is implemented in C++20 with OpenMP and compiled with \verGcc{} at
optimisation level \texttt{-O3}. It provides a verification mode that compares the
programme with exhaustive enumeration, a verification mode for the vertex theorem, and a
front mode that writes every Pareto point together with its decisions.

%=====================================================================================
\section{Experimental design}\label{sec:design}

\paragraph{Instances}
We use the two items of the industrial case of \citet{Kania2022CIE}. Item~1 has unit
price $p=134$, lead time $L=5$, $\moq=r=70$, mean weekly demand $\mu=163$ and standard
deviation $\sigma=41$, with $\MS=1.0$. Item~2 has $p=483.85$, $L=4$, $\moq=r=3$,
$\mu=12.2$, $\sigma=3.6$ and $\MS=1.4$. Both have ordering cost $c=200$ and holding cost
rate $h=0.1p/52$ per week. These values were reconstructed from the published description
and the authors' open-source code. Demand sample paths are drawn from a normal
distribution with the given moments, rounded and truncated at zero, for seeds 1--10. The
initial stock covers the demand during the lead time plus one mean period. Grid mode uses
the grid of the original study; continuous mode uses the full domain $\Omega$.

\paragraph{Experiments}
E1 checks correctness: the programme is compared with exhaustive enumeration on 100
random small instances, and the vertex theorem is checked on $10^6$ random points of
$\Omega$ for each of 20 instances. E2 computes the complete fronts at $T=48$ for both
items, both outer domains and ten demand seeds. E3 measures scaling: one inner problem
for horizons $T=12,\dots,96$ (three seeds), and the complete continuous front of Item~1
at $T=96$. E4 varies the demand coefficient of variation, the rounding value and the
ordering cost. E5a compares the inner programme on the cost--holding pair with an
$\varepsilon$-constraint mixed-integer programme \citep{Haimes1971} solved with HiGHS
\citep{HuangfuHall2018} through SciPy \verScipy, at one outer point and growing horizon.
E5b runs NSGA-III \citep{DebJain2014} in pymoo \verPymoo{} \citep{BlankDeb2020} with 455
Das--Dennis reference directions for one hour per run, with and without a problem-specific
repair operator that enforces feasibility, the ordering rule of \cref{lem:D} and the
lattice of order quantities, and compares the archives with the exact fronts.
E6 solves the public benchmark described next.

\paragraph{Public benchmark (E6)}
The two industrial items and their normally distributed demand do not show how the
algorithm behaves on the irregular demand of practice. We therefore built a benchmark from
two public sources: the M5 data of Walmart unit sales at the item--store level
\citep{Makridakis2022M5}, and the UCI Online Retail II transactions of a British online
retailer \citep{Chen2012Retail}, which are licensed under CC BY 4.0. Sales are aggregated to
weeks; for Online Retail II, cancelled quantities are netted against the original sales,
so that orders cancelled at once do not appear as demand. Each planning window of $T=48$ weeks is preceded by 52 weeks of history, and all
forecast and cost parameters are estimated from the history only, so that no information
from the planning window is used. The mean $\mu$ and standard deviation $\sigma$ (at least
$0.5$) of weekly demand serve as forecast parameters, and the unit price $p$ is the median
selling price. Series are classified as smooth, erratic, intermittent or lumpy by the
average inter-demand interval and the squared coefficient of variation of non-zero demand
\citep{SyntetosBoylan2005}, and a stratified random sample is drawn from each class. M5
contributes the last three non-overlapping windows of each selected series and Online
Retail II the last window. Item parameters follow rules fixed before any run: lead time
$L=2$ weeks, case pack $r=\max\{1,\operatorname{round}(\mu/2)\}$ with $\moq=r$ (so that
\cref{as:lattice} holds), holding cost $h=0.1p/52$, $\MS=1$, and an ordering cost chosen so
that the classical economic order quantity covers four weeks of mean demand. Series with
fewer than four weeks of non-zero demand in any history, or without a price, are excluded.
Every run is limited to \EsixTlim{} minutes of wall-clock time on six threads; a run that
reaches the limit is reported as unsolved. To verify exactness on real demand, 30 randomly
chosen instances are also solved without the pruning rules. The M5 data cannot be redistributed, so the repository contains the
script that rebuilds the instances from the official files; the Online Retail II instances
are included directly.

\paragraph{Quality indicators}
For approximate sets we report the share of points that lie on the exact front, the
share dominated by it, the share of the exact front recovered, the IGD$^+$ indicator
\citep{Ishibuchi2015} and the additive $\varepsilon^+$ indicator \citep{Zitzler2003}, all
on objectives normalised by the range of the exact front. Equality uses absolute
tolerances of $10^{-4}$, $10^{-6}$, $10^{-8}$ and $10^{-8}$ for the four objectives, just
above the precision of the output files.

\paragraph{Computing environment}
All experiments ran on one laptop (\verCpu, 12 logical cores, 14\,GiB memory) under
\verOs, using \verGcc, CMake \verCmake{} and Ninja \verNinja. Analyses and figures were produced with
Python \verPython, NumPy \verNumpy{} and Matplotlib \verMatplotlib. Every table and figure
in this paper is generated by one script from the raw output files.

%=====================================================================================
\section{Results}\label{sec:results}

\paragraph{Correctness}
The programme agreed with exhaustive enumeration on all 100 small instances. For every
one of the $2\times10^7$ random points used to test \cref{thm:vertex}, a dominating vertex
was found. The mixed-integer baseline of E5a produced exactly the same fronts as the
programme at every horizon at which it finished (\cref{tab:e5a}). Over all runs of E5b,
not a single exact point was dominated by any of the roughly $1.2\times10^8$ solutions
evaluated by NSGA-III. During the formal analysis we found that an earlier version of the
ordering rule omitted the second affected period. It removed 2 of 19\,372 Pareto points in
one instance (Item~2, seed~42). A second, purely numerical defect was found when the
faster algorithm was compared with the original one: the final filter compared service
levels as floating-point values with a small tolerance, and near $\CSL=1$, where
$\Phi$ is not resolved in double precision, this occasionally discarded non-dominated
points. The service level is now compared through $B=\Ss+\mu\SOT$, which is exact because
$\Phi$ is strictly increasing. All fronts reported here were produced with both
corrections, and the complete fronts of E2, E4 and of 30 randomly chosen E6 instances are
identical, within the printed precision, to those of the period-by-period programme run
over all classes and without any pruning.

\paragraph{Complete fronts at $T=48$}
\Cref{tab:e2} summarises E2. A complete front contains tens of thousands of
Pareto optimal plans and is computed in a few seconds. Continuous safety parameters
enlarge the front by roughly one fifth for Item~1 and double it for Item~2. The
computing time grows by about 40\% for Item~1 but fivefold for Item~2, because the small
rounding value of Item~2 creates many requirement lines and hence many arrangement
vertices and classes. \Cref{fig:front} shows two
projections of one front. The cost--holding projection is split into separate branches,
one for each total order quantity. The holding--turnover projection shows how a higher
service level shifts the attainable turnover downwards.

\begin{table}[t]
  \centering
  \caption{Complete Pareto fronts at $T=48$ (E2, mean $\pm$ standard deviation over ten
  demand seeds). Candidates are grid points or arrangement vertices; classes are distinct
  rounded requirement vectors, that is, inner problems solved.}
  \label{tab:e2}
  \small
  \tabin{tab_e2.tex}
\end{table}

\begin{figure}[t]
  \centering
  \figin{\textwidth}{fig_front.pdf}
  \caption{Exact Pareto front of Item~1, seed~1, continuous safety parameters, coloured by
  cycle service level: (a) purchasing and ordering cost against holding cost; (b) holding
  cost against inventory turnover.}
  \label{fig:front}
\end{figure}

\paragraph{Scaling}
\Cref{fig:scaling} and \cref{tab:e3fit} report the growth of one inner problem with the
horizon. The fitted exponents of time against $T$ are 5.94 and 6.53, and those of the
labels generated are 4.88 and 5.04. Both agree with the analysis of \cref{app:complexity}:
with $S=\Theta(T)$ states per layer and $O(T^2S)$ labels per node, the total number of
labels is $O(T\cdot S\cdot T^2S)=O(T^5)$ and the time is $O(T^6\log T)$ for a fixed ratio
$\mu/r$. The slightly larger time exponent of Item~2 is consistent with the logarithmic
factor and with pre-asymptotic effects at these horizons. The bound thus describes the
observed growth, not only the worst case. The complete continuous front of
Item~1 at $T=96$, with \EfullFront{} Pareto points, was computed in \EfullTime{} minutes.

\begin{figure}[t]
  \centering
  \figin{\textwidth}{fig_scaling.pdf}
  \caption{One inner dynamic programme against the horizon $T$ (E3, mean over three
  seeds): (a) computing time; (b) labels generated. Slopes are least-squares fits on
  logarithmic axes for $T\ge 36$.}
  \label{fig:scaling}
\end{figure}

\begin{table}[t]
  \centering
  \caption{Fitted scaling exponents of one inner problem (E3).}
  \label{tab:e3fit}
  \small
  \tabin{tab_e3fit.tex}
\end{table}

\paragraph{Sensitivity}
\Cref{tab:e4} reports E4. A higher demand variability multiplies the number of classes
and increases the computing time, but the front becomes smaller, from about 48 thousand
to 28 thousand points, because fewer plans remain non-dominated when requirements vary
strongly between periods. A smaller rounding value refines the lattice of order
quantities and enlarges both the front and the computing time considerably ($r=35$:
135 thousand points in 71 seconds), whereas $r=140$ reduces them to 15 thousand points in
under one second. The ordering cost leaves the number of classes, the front size and the
computing time unchanged. In
\EinvOK{} of \EinvN{} seeds the fronts for $c=50$, $200$ and $800$ consisted of the same
plans, compared through the cost-independent signature of each point (total quantity,
number of orders, holding cost, service level and turnover). \Cref{prop:c} guarantees this for $c<218.1$; the identity at $c=800$
shows that the condition of the proposition is sufficient but not necessary.

\begin{table}[t]
  \centering
  \caption{Sensitivity for Item~1, $T=48$, continuous domain (E4, five seeds).}
  \label{tab:e4}
  \small
  \tabin{tab_e4.tex}
\end{table}

\paragraph{Comparison with an exact generic method}
The $\varepsilon$-constraint programme solves two mixed-integer problems (one per
lexicographic stage) for every admissible number of orders, and its time grows steeply
with the horizon (\cref{tab:e5a,fig:milp}). Both methods return identical fronts at every
horizon. At $T=48$ a single outer point takes 1\,262 and 1\,535 seconds for the two items,
against 18 and 31 milliseconds for the dynamic programme, a ratio of about $5$--$7\times10^4$.
The ratio itself grows with the horizon, and a complete front requires several hundred
such outer points.

\begin{table}[t]
  \centering
  \caption{Dynamic programme against the $\varepsilon$-constraint MILP on the bi-objective
  (POC, HC) inner problem at $\Ss=8$, $\SOT=0.3$ (E5a, seed 1).}
  \label{tab:e5a}
  \small
  \tabin{tab_e5a.tex}
\end{table}

\begin{figure}[t]
  \centering
  \figin{0.6\textwidth}{fig_milp.pdf}
  \caption{Computing time of the dynamic programme and of the $\varepsilon$-constraint MILP
  against the horizon (E5a). Dynamic-programme times are measured with a resolution of one
  millisecond.}
  \label{fig:milp}
\end{figure}

\paragraph{Comparison with NSGA-III}
\Cref{tab:e5b} and \cref{fig:nsga} show that one hour of NSGA-III, more than two orders
of magnitude longer than the exact computation, yields sets that are almost entirely
dominated. Without repair, no archive point lies on the exact front. The repair operator
enlarges the archives about fivefold and reduces the $\varepsilon^+$ indicator by roughly
one order of magnitude (for Item~1 against the continuous front from $0.277$ to $0.034$,
for Item~2 from $0.474$ to $0.039$). Even so, 97.8--99.4\% of the archive points remain
dominated, and on average only 0.14--0.63\% of the exact front is recovered. Indicator values that look acceptable in isolation thus hide a
large gap, which becomes visible only against the exact reference.

\begin{table}[t]
  \centering
  \caption{NSGA-III archives after one hour against the exact fronts (E5b, seeds 1--3).
  $|A|$ is the archive size; the last column counts exact points dominated by an archive
  point and must be zero.}
  \label{tab:e5b}
  \scriptsize
  \resizebox{\textwidth}{!}{\tabin{tab_e5b.tex}}
\end{table}

\begin{figure}[t]
  \centering
  \figin{0.6\textwidth}{fig_nsga.pdf}
  \caption{Item~1, seed~1: the (POC, HC) frontier of the exact front and the archives of
  NSGA-III without and with repair (E5b).}
  \label{fig:nsga}
\end{figure}

\paragraph{Public benchmark}
\Cref{tab:e6} and \cref{fig:e6} report E6. Of the \EsixInst{} instances (\EsixSeries{}
series), \EsixSolved{} (\EsixPctSolved\%) were solved exactly within the time limit, with a
median time of \EsixMedTime{} seconds and a maximum of \EsixMaxTime{} seconds; the solved
instances required \EsixTotalMin{} minutes of computing time in total. The remaining
\EsixTimeout{} instances reached the limit. They are explained by the complexity analysis
rather than by the demand class. The number of states per layer grows with the realised
demand of the planning window divided by the case pack, $S=\Theta(\sum_t D(t)/r)$, whereas
the case pack is set from the history. Series whose sales rise sharply between history and
planning window, typically products that were introduced or out of stock during the
history, therefore yield a ratio of window demand to case pack far above the intended
value of about two. \Cref{fig:e6}(c) shows that computing time is governed by this ratio
across all classes, as predicted by the bound of \cref{app:complexity}. In \EsixChkOK{} of \EsixChkN{}
reference runs without pruning, the front was identical to that of the algorithm, which
confirms the reductions of \cref{sec:method} on irregular real demand.

\begin{table}[t]
  \centering
  \caption{Public benchmark (E6): median and interquartile range per source and demand
  class. ``At limit'' counts runs stopped by the time limit; $\mu$, $\mu/r$ refer to all
  instances, the remaining columns to solved instances. Times are per instance with the
  continuous outer domain.}
  \label{tab:e6}
  \resizebox{\textwidth}{!}{\tabin{tab_e6.tex}}
\end{table}

\begin{figure}[t]
  \centering
  \figin{\textwidth}{fig_e6.pdf}
  \caption{Public benchmark (E6): (a) demand classes of the selected series by average
  inter-demand interval (ADI) and squared coefficient of variation; dashed lines are the
  cut-offs of \citet{SyntetosBoylan2005}; circles M5, triangles Online Retail II;
  (b) computing time per solved instance by class (S smooth, E erratic, I intermittent,
  L lumpy); (c) computing time against the realised weekly demand of the planning window
  divided by the case pack $r$, with runs stopped at the time limit marked by crosses.}
  \label{fig:e6}
\end{figure}

%=====================================================================================
\section{Discussion}\label{sec:discussion}

The first implication concerns interactive decision support. Methods such as E-NAUTILUS
are designed to navigate over a precomputed set of solutions. Our results show that for
the model of \citet{Kania2022CIE} this set can be the complete Pareto front, computed in
seconds, rather than an approximation of unknown quality. Every plan the decision maker
sees is then Pareto optimal, and every attainable trade-off is available. The
comparison with NSGA-III shows that this difference matters in practice: the
approximation offered plans that were almost all dominated, and its standard indicators
did not reveal it.

The second implication is methodological. The exactness of the algorithm rests on
reductions whose safety is easy to assume and hard to verify. Our own first version of the
ordering rule overlooked a feasibility condition that affects only plans with very low
stock. It passed exhaustive tests on small instances and was found only through the
formal proof. We therefore regard the combination of proofs and three independent
numerical checks (enumeration, an independent mixed-integer model, and dominance against
very large approximate sets) as necessary rather than optional for exact multiobjective
algorithms.

The study has limitations. The model is single-item and evaluates each plan on one
demand sample path, as in the original study, so the front describes trade-offs for a
given forecast rather than in expectation. The running time grows as $T^6$, so very long
horizons or very small rounding values will need approximation with guarantees; an
$\varepsilon$-Pareto variant of the label-setting procedure is a natural candidate. The
service-level objective saturates near one, and many Pareto points differ only in its
eighth decimal; a decision maker would treat these as equivalent, which suggests
presenting the front with a resolution chosen by the user. The cycle service level of the model assumes normally distributed demand, which fits
intermittent and lumpy series poorly; on these classes the service objective should be
read as a nominal index rather than as a probability. Finally, the parameters of the
industrial items were reconstructed from published information and may differ in detail
from the company data.

These limitations point to the next steps. For many items, the exact single-item fronts
can replace the representative-item approximation of DESMILS
\citep{Kania2024DESMILS}. For larger instances, preference-conditioned learning models
such as Pareto set learning \citep{Lin2022PSL} can be trained and, for the first time
in this application, evaluated against exact fronts rather than against other
heuristics.

%=====================================================================================
\section{Conclusion}\label{sec:conclusion}

We have shown that the four-objective lot-sizing model with safety stock and safety order
time of \citet{Kania2022CIE} admits an exact and fast solution method. A cumulative
reformulation turns the inner problem into a multi-criteria shortest path problem. A
label-setting dynamic programme with three proven reductions solves it, and a vertex
theorem reduces continuous safety parameters to a finite candidate set. On industrial
instances the complete fronts, with tens of thousands of Pareto optimal plans, are
obtained in seconds. This is more than four orders of magnitude faster than an exact
mixed-integer approach, and it far exceeds the quality of a state-of-the-art
evolutionary algorithm. Exact fronts make interactive lot sizing verifiable and provide
the reference that learning-based methods need.

\section*{Data and code availability}
The C++ implementation, the baseline scripts, the instance generator, the benchmark
instances with their exact Pareto fronts, and the script that produces every table and
figure are available at \url{https://github.com/USERNAME/lsdp}. The M5 data are available
from the competition website and are not redistributed; the repository contains the script
that rebuilds the M5 instances from the official files. The UCI Online Retail II data are
used under the CC BY 4.0 licence.

\section*{Declaration of generative AI use}
During the preparation of this work the author used an AI assistant to help draft
code and text. The author reviewed and edited all content and takes full responsibility
for the publication.

\appendix
\renewcommand{\thesection}{\Alph{section}}
\section{Theoretical foundations}\label{app:theory}
% Body of the theoretical appendix; \input by main.tex.
\noindent
This appendix states the model in the form used by the algorithm, proves that the
algorithm returns the exact Pareto front, and derives its complexity. Every result
is linked to the experiment that checks it numerically (\cref{app:verification}).
Throughout, \emph{dominance} means Pareto dominance: an outcome is weakly better in
every objective and strictly better in at least one. Two plans with identical
objective vectors are treated as one point of the front.

%=====================================================================================
\subsection{Model and notation}\label{app:model}

The planning horizon has periods $t=1,\dots,T$ and a deterministic lead time $L$, so
orders can be placed in the decision periods $\tau=1,\dots,\Tp$ with $\Tp=T-L$; an order
placed in period $\tau$ is received in period $\tau+L$. Demand $D(t)\ge 0$ is given
(one sample path of the stochastic demand), $\cumD(t)=\sum_{s\le t}D(s)$ with
$\cumD(0)=0$, and $I_0$ is the initial stock. Order quantities lie in
\begin{equation}
  \cQ=\{0\}\cup\{\moq+a\,r : a\in\mathbb{N}_0\},
\end{equation}
where $\moq$ is the minimum order quantity and $r$ the rounding value. The outer
(policy) variables are the safety stock $\Ss$ and the safety order time $\SOT$, taken
from the domain
\begin{equation}
  \Omega=\{(\Ss,\SOT)\in\mathbb{R}^2_{\ge 0} : \SOT+\Ss/\mu\le\MS\},
\end{equation}
or from a finite grid in $\Omega$ (\emph{grid mode}).

\paragraph{Cumulative reformulation.}
Let $X(\tau)=\sum_{s\le\tau}Q(s)$, $X(0)=0$. Then the stock at the end of period
$\tau+L$ is $I(\tau+L)=I_0+X(\tau)-\cumD(\tau+L)$, while $I(t)=I_0-\cumD(t)$ for
$t\le L$ does not depend on the decisions. The service requirement
$I(t)\ge\Ss+\SOT\cdot D(t)$ for $t=L+1,\dots,T$ becomes
\begin{equation}\label{eq:req}
  X(\tau)\ge\beta_\tau(\Ss,\SOT)
  :=\Ss+\SOT\cdot D(\tau+L)+\cumD(\tau+L)-I_0,\qquad \tau=1,\dots,\Tp,
\end{equation}
and the average stock $\Ibar(t)=\tfrac12\bigl(I(t-1)+I(t)\bigr)$ in period $t=\tau+L$ is
\begin{equation}\label{eq:ibar}
  \Ibar(\tau+L)=b_\tau+\tfrac12\bigl(X(\tau-1)+X(\tau)\bigr),\qquad
  b_\tau:=I_0-\tfrac12\bigl(\cumD(\tau+L-1)+\cumD(\tau+L)\bigr).
\end{equation}
Inventory turnover is defined only for positive average stock, so we impose
$\Ibar(t)\ge\eps$ for all $t$, with a fixed threshold $\eps>0$ (the implementation
uses $\eps=1$ unit).

\paragraph{Objectives.}
With $k=\#\{\tau : Q(\tau)>0\}$ the number of orders, the four objectives of
Kania et al.\ read
\begin{align}
  \POC &= p\,X(\Tp)+c\,k &&\text{(min)},\label{eq:poc}\\
  \HC  &= h\sum_{t=1}^{T}\Ibar(t)
        = h\Bigl[C_0+\sum_{\tau=1}^{\Tp}X(\tau)-\tfrac12X(\Tp)\Bigr]
        &&\text{(min)},\label{eq:hc}\\
  \CSL &= \Phi\bigl(B/\sigma\bigr),\quad B:=\Ss+\mu\,\SOT &&\text{(max)},\label{eq:csl}\\
  \ITO &= A_0+\sum_{\tau=1}^{\Tp}\frac{D(\tau+L)+\sigma}{\Ibar(\tau+L)} &&\text{(max)},\label{eq:ito}
\end{align}
where $C_0=\sum_{t\le L}\Ibar(t)+\sum_{\tau}b_\tau$ and
$A_0=\sum_{t\le L}(D(t)+\sigma)/\Ibar(t)$ are constants. The second form of
\eqref{eq:hc} follows from \eqref{eq:ibar} and $X(0)=0$. We write
$F=(\POC,\HC,-\CSL,-\ITO)$, to be minimised componentwise.

\paragraph{Assumptions.} Both instances studied satisfy the following.
\begin{assumption}\label{as:initial}
  $I(L)=I_0-\cumD(L)\ge 0$ and $\Ibar(t)\ge\eps$ for $t\le L$.
\end{assumption}
\begin{assumption}\label{as:lattice}
  $\moq$ is a positive multiple of $r$, i.e.\ $\moq=m r$ with $m\in\mathbb{N}$.
\end{assumption}
Under \cref{as:lattice}, $\cQ\setminus\{0\}=\{r j : j\ge m\}$ is closed under addition,
and the set of reachable cumulative quantities is
$\cL=\{0\}\cup\{r j : j\ge m\}\subset r\mathbb{N}_0$. For $y\in\mathbb{R}$ let
$\lceilL{y}=\min\{\ell\in\cL : \ell\ge y\}$. Because $X(\tau)\in\cL$,
\begin{equation}\label{eq:kappa}
  X(\tau)\ge\beta_\tau \iff X(\tau)\ge\kappa_\tau,\qquad
  \kappa_\tau:=\lceilL{\beta_\tau}.
\end{equation}
We call $\kappa=(\kappa_1,\dots,\kappa_{\Tp})$ the \emph{rounded requirement vector}.

%=====================================================================================
\subsection{The inner problem as a multi-criteria shortest path}\label{app:inner}

Fix an outer point $z=(\Ss,\SOT)$. Then $\CSL$ is constant, and a plan is a
non-decreasing path $X(0)=0,X(1),\dots,X(\Tp)$ in $\cL$ whose increments lie in $\cQ$.
Consider the layered directed acyclic graph whose nodes are pairs $(\tau,x)$ with
$x\in\cL$, and whose arcs $(\tau-1,x')\to(\tau,x)$ exist when
\begin{equation}\label{eq:arc}
  x-x'\in\cQ,\qquad x\ge\kappa_\tau,\qquad b_\tau+\tfrac12(x'+x)\ge\eps .
\end{equation}
Each arc carries the additive cost vector
\begin{equation}\label{eq:arccost}
  \bigl(\,\mathbb{1}[x>x'],\; x,\; \tfrac{D(\tau+L)+\sigma}{b_\tau+\frac12(x'+x)}\,\bigr),
\end{equation}
and a path to $(\tau,x)$ accumulates the label $\ell=(k,W,R)$: the number of orders, the
partial sum $W=\sum_{s\le\tau}X(s)$, and the partial turnover $R$. At a terminal node
$(\Tp,x)$, \eqref{eq:poc}--\eqref{eq:ito} give
\begin{equation}\label{eq:terminal}
  \POC=p x+c k,\qquad \HC=h\bigl(C_0+W-\tfrac12x\bigr),\qquad \ITO=A_0+R .
\end{equation}
For a fixed terminal node these are strictly monotone in the label (increasing in $k$
and $W$, increasing in $R$ for the maximised $\ITO$). Hence two paths ending at the same
node compare in $F$ exactly as their labels compare under the order
\begin{equation}
  (k,W,R)\preceq(k',W',R')\iff k\le k',\ W\le W',\ R\ge R'.
\end{equation}

The label-setting algorithm processes the layers $\tau=1,\dots,\Tp$ in order, extends
every stored label along every outgoing arc, and keeps at each node only the labels that
are non-dominated under $\preceq$. The implementation stores them as one
two-dimensional staircase in $(W,R)$ per value of $k$; this is a data structure for the
same dominance test and does not affect correctness.

%=====================================================================================
\subsection{Reductions of the search space}\label{app:reductions}

The algorithm searches a restricted graph: states are bounded above (\cref{lem:ub}), and
ordering arcs that can never be Pareto-relevant are removed (\cref{lem:D}). Both lemmas
are proved by exhibiting, for every plan that the restriction excludes, a feasible plan
with the same outer point that dominates it. So no Pareto-optimal plan is excluded.

\begin{lemma}[Upper bound on the cumulative quantity]\label{lem:ub}
  Let $\hat\beta=\max\{0,\max_\tau\beta_\tau\}$ and let \cref{as:initial} hold. Every
  feasible plan with $X(\Tp)\ge\hat\beta+\max\{\moq,r\}+2\eps$ is dominated by a feasible
  plan with the same outer point. Consequently every Pareto-optimal plan satisfies
  $X(\Tp)<\hat\beta+\max\{\moq,r\}+2\eps$.
\end{lemma}

\begin{proof}
  Let $\tau^\ast$ be the last period with $Q(\tau^\ast)>0$ (it exists because
  $X(\Tp)>0$), so $X(\tau)=X(\Tp)$ for $\tau\ge\tau^\ast$. Reduce the last order by
  $\delta=r$ if $Q(\tau^\ast)\ge\moq+r$, and by $\delta=\moq$ (that is, remove it) if
  $Q(\tau^\ast)=\moq$. In both cases $Q(\tau^\ast)-\delta\in\cQ$, $\delta\le\max\{\moq,r\}$,
  and the new plan $X'$ equals $X$ before $\tau^\ast$ and $X-\delta$ from $\tau^\ast$ on;
  it remains non-decreasing.

  \emph{Feasibility.} For $\tau\ge\tau^\ast$,
  $X'(\tau)=X(\Tp)-\delta\ge\hat\beta+2\eps\ge\beta_\tau+2\eps$, so \eqref{eq:req} holds and
  $I'(\tau+L)=X'(\tau)-\beta_\tau+\Ss+\SOT\,D(\tau+L)\ge 2\eps$. Hence
  $\Ibar'(t)\ge2\eps$ for $t>\tau^\ast+L$. In period $t=\tau^\ast+L$,
  $\Ibar'(t)=\tfrac12\bigl(I(t-1)+I'(t)\bigr)\ge\tfrac12(0+2\eps)=\eps$, because $I(t-1)$
  is unchanged and non-negative: by \cref{as:initial} if $\tau^\ast=1$, and by
  \eqref{eq:req} in period $\tau^\ast-1$ otherwise. Earlier periods are unchanged.

  \emph{Dominance.} $\POC$ decreases by $p\delta>0$ (and by a further $c$ if the order is
  removed). By \eqref{eq:hc}, $\HC$ changes by
  $h\bigl[-\delta(\Tp-\tau^\ast+1)+\delta/2\bigr]<0$. $\CSL$ is unchanged. Every
  $\Ibar(t)$ either stays the same or decreases, so $\ITO$ does not decrease.
\end{proof}

The implementation therefore uses the state set
$\{x\in\cL : x\le\lceil\hat\beta\rceil+\max\{\moq,r\}+\lceil2\eps\rceil\}$, where
$\hat\beta$ is the maximum of $\beta_\tau$ over the whole outer domain. This is one
superset that serves every outer point.

\begin{lemma}[Order only when needed]\label{lem:D}
  Let \cref{as:lattice} hold, and consider a feasible plan with an order $q=Q(\tau)>0$ and
  $x'=X(\tau-1)$. Suppose that
  \begin{enumerate}[label=(\alph*),nosep]
    \item $x'\ge\beta_\tau$ (equivalently $x'\ge\kappa_\tau$),
    \item $b_\tau+x'\ge\eps$, and
    \item $\tau=\Tp$, or $b_{\tau+1}+x'+q/2\ge\eps$.
  \end{enumerate}
  Then the plan is dominated by a feasible plan with the same outer point.
\end{lemma}

\begin{proof}
  If $\tau<\Tp$, move the order to the next period:
  $Q'(\tau)=0$ and $Q'(\tau+1)=Q(\tau+1)+q$. This quantity lies in $\cQ$, either because
  $Q(\tau+1)=0$ or because $\cQ\setminus\{0\}$ is closed under addition
  (\cref{as:lattice}). If $\tau=\Tp$, drop the order. The new plan differs from the old
  one only in $X'(\tau)=x'$.

  \emph{Feasibility.} The requirement in period $\tau$ holds by (a), and all other
  requirements are unchanged. By \eqref{eq:ibar} and (b),
  $\Ibar'(\tau+L)=b_\tau+x'\ge\eps$. If $\tau<\Tp$, then $X(\tau+1)\ge X(\tau)=x'+q$, and (c)
  gives $\Ibar'(\tau+L+1)=b_{\tau+1}+\tfrac12\bigl(x'+X(\tau+1)\bigr)\ge b_{\tau+1}+x'+q/2\ge\eps$.
  No other period changes.

  \emph{Dominance.} If $\tau<\Tp$, then $X(\Tp)$ is unchanged and $k$ stays the same or
  decreases by one, so $\POC$ does not increase. $\sum_\tau X(\tau)$ decreases by $q$, so
  $\HC$ decreases strictly. If $\tau=\Tp$, then $\POC$ decreases and $\HC$ changes by
  $h(-q+q/2)<0$. $\CSL$ is unchanged, and $\ITO$ does not decrease because the average
  stock in periods $\tau+L$ and $\tau+L+1$ does not increase.
\end{proof}

\begin{remark}[Why condition (c) is needed]\label{rem:c}
  Postponing an order lowers the average stock in \emph{two} periods, $\tau+L$ and
  $\tau+L+1$. A version of \cref{lem:D} that checks only (a) and (b) is not valid: the
  postponed plan may violate $\Ibar\ge\eps$ in period $\tau+L+1$, so the claimed dominating
  plan need not exist. In our experiments the earlier implementation, which omitted (c),
  missed 2 of 19\,372 Pareto points in one instance (SKU~2, seed~42, grid mode; both
  points have $\Ss=\SOT=0$, where the stock is lowest). The corrected rule reproduces the
  front of the DP without any pruning. All 20 instances of experiment E2 (seeds 1--10) are
  unaffected.
\end{remark}

\begin{remark}[Implementation]
  For fixed $(\tau,x')$, condition (c) is monotone in $q$, and (a) and (b) do not depend
  on $q$. Once an order size is prunable, every larger one is too. The implementation
  enumerates target states in increasing order and stops at the first prunable order.
\end{remark}

%=====================================================================================
\subsection{Exactness of the algorithm}\label{app:exact}

\begin{theorem}[Exactness for a fixed outer point]\label{thm:exact}
  Let \cref{as:initial,as:lattice} hold and fix $z\in\Omega$. Let $G_z$ be the graph of
  \cref{app:inner} restricted to the state set of \cref{lem:ub}, with the arcs excluded by
  \cref{lem:D} removed. The label-setting algorithm on $G_z$, followed by evaluating
  \eqref{eq:terminal} and filtering non-dominated outcomes over all terminal nodes,
  returns exactly the set of non-dominated outcome vectors of the inner problem at $z$.
\end{theorem}

\begin{proof}
  \emph{(i) The restriction loses no Pareto-optimal plan.} A Pareto-optimal plan is not
  dominated by any feasible plan. By \cref{lem:ub} it lies within the state bound, and by
  \cref{lem:D} none of its arcs is pruned. So it is a path of $G_z$. Conversely, every
  path of $G_z$ satisfies \eqref{eq:arc} and is therefore a feasible plan.

  \emph{(ii) Label dominance is preserved by extension.} Arc existence and arc cost
  depend only on the arc $(\tau-1,x')\to(\tau,x)$, not on the path that reached
  $(\tau-1,x')$. This holds for the pruning rule as well, which depends on
  $(\tau,x',x-x')$ only. If labels $\ell\preceq\ell'$ sit at the same node, extending both
  along the same arc preserves $\preceq$, because the costs \eqref{eq:arccost} add
  componentwise.

  \emph{(iii) Induction over layers.} Suppose that at layer $\tau-1$ every path of $G_z$
  to a node has a stored label at that node which weakly dominates it under $\preceq$.
  Extending that stored label along the path's next arc weakly dominates the extended
  path, by (ii). Either it is kept, or it is discarded because another kept label
  dominates it; by transitivity, some kept label dominates the extended path. So the claim
  holds at layer $\tau$, and it holds trivially at layer $0$.

  \emph{(iv) Terminal evaluation.} By (iii) and the monotonicity of \eqref{eq:terminal},
  every plan in $G_z$ ending at $(\Tp,x)$ is weakly dominated in $F$ by a returned outcome
  at the same node. Combining this with (i), the returned set contains every
  non-dominated outcome. It contains only feasible outcomes, so after the final filter it
  equals the set of non-dominated outcome vectors.
\end{proof}

\begin{lemma}[The inner problem depends on $z$ only through $\kappa$]\label{lem:memo}
  For $z,z'\in\Omega$ with $\kappa(z)=\kappa(z')$, the graphs $G_z$ and $G_{z'}$ coincide,
  so the inner problems have the same outcomes in $(\POC,\HC,\ITO)$. If
  $\kappa(z)\le\kappa(z')$ componentwise, every plan feasible at $z'$ is feasible at $z$.
\end{lemma}

\begin{proof}
  By \eqref{eq:kappa}, $z$ enters \eqref{eq:arc} only through $\kappa_\tau$. Condition
  (a) of \cref{lem:D} is equivalent to $x'\ge\kappa_\tau$, and conditions (b) and (c) do
  not involve $z$. The objectives \eqref{eq:terminal} do not involve $z$, and the state
  bound is common to all outer points. The second statement follows because
  $x\ge\kappa_\tau(z')\ge\kappa_\tau(z)$.
\end{proof}

\begin{corollary}[Exact memoization]\label{cor:memo}
  Among outer points with the same rounded requirement vector, only one with the largest
  $B=\Ss+\mu\SOT$ can contribute to the Pareto front: the others have identical
  $(\POC,\HC,\ITO)$ outcomes and a strictly smaller $\CSL$, because $\Phi$ is strictly
  increasing. The algorithm therefore solves one inner problem per distinct vector
  $\kappa$, attaches the largest $B$ of its class, and applies a final four-objective
  filter over the classes. By \cref{thm:exact} the result is the exact front over the
  outer candidate set.
\end{corollary}

%=====================================================================================
\subsection{The continuous outer domain}\label{app:continuous}

Each $\beta_\tau$ is affine in $z$: $\beta_\tau(z)=\Ss+D_\tau\SOT+c_\tau$ with
$D_\tau=D(\tau+L)$ and $c_\tau=\cumD(\tau+L)-I_0$. Hence $\kappa_\tau(z)$ is piecewise
constant. It can change only on the lines
\begin{equation}
  H_{\tau,\lambda}=\{z : \Ss+D_\tau\SOT=\lambda-c_\tau\},\qquad
  \lambda\in\Lambda:=\cL\cap\bigl[0,\lceilL{\hat\beta}\bigr],
\end{equation}
where $\hat\beta$ is the maximum of $\beta_\tau$ over $\Omega$. Let $\mathcal{A}$ be the
arrangement formed by these lines together with the three lines that bound $\Omega$
($\Ss=0$, $\SOT=0$, $\SOT+\Ss/\mu=\MS$). Let $V$ be the finite set of its vertices,
that is, the intersection points of two non-parallel lines of $\mathcal{A}$ that lie in
$\Omega$.

\begin{theorem}[Vertex sufficiency]\label{thm:vertex}
  For every $z\in\Omega$ there is a vertex $v\in V$ with $\kappa(v)\le\kappa(z)$
  componentwise and $B(v)\ge B(z)$. Consequently the set of non-dominated outcome vectors
  over all of $\Omega$ equals the set of non-dominated outcome vectors over $V$.
\end{theorem}

\begin{proof}
  Let $P(z)=\{w\in\Omega : \beta_\tau(w)\le\kappa_\tau(z)\ \forall\tau\}$. It is the
  intersection of the compact polygon $\Omega$ with finitely many closed half-planes. It
  is therefore a compact convex polygon, and it contains $z$ because
  $\beta_\tau(z)\le\kappa_\tau(z)$. The linear function $B$ attains its maximum over
  $P(z)$ at a vertex $v$ of $P(z)$, so $B(v)\ge B(z)$. Every $w\in P(z)$ satisfies
  $\kappa_\tau(w)\le\kappa_\tau(z)$, because $\kappa_\tau(z)\in\cL$ is an upper bound of
  $\beta_\tau(w)$ and $\kappa_\tau(w)$ is the smallest such element of $\cL$. A vertex of
  $P(z)$ is the intersection of two non-parallel boundary lines of $P(z)$. Each such line
  is either a boundary line of $\Omega$ or a line $\beta_\tau=\kappa_\tau(z)$, that is,
  $H_{\tau,\lambda}$ with $\lambda=\kappa_\tau(z)\in\Lambda$. Hence $v\in V$.

  For the consequence, let $y=F(z,X)$ be non-dominated over $\Omega$. By
  \cref{lem:memo}, the plan $X$ is feasible at $v$, and $F(v,X)$ has the same
  $(\POC,\HC,\ITO)$ and $\CSL(v)\ge\CSL(z)$. So $F(v,X)$ weakly dominates $y$. Since $y$
  is non-dominated, $F(v,X)=y$; thus $y$ is attained at a vertex. The reverse inclusion
  holds because $V\subseteq\Omega$.
\end{proof}

\begin{remark}[Correction of an earlier claim]
  An earlier version of this work asserted that the continuous front need not be closed,
  so that some cells of the arrangement contributed points only as a supremum of
  $\CSL$. To catch these, it added one interior probe per angular sector around each
  vertex. \Cref{thm:vertex} shows the probes are unnecessary. The vertex of a cell where
  the supremum is attained belongs to a neighbouring cell with a componentwise smaller
  requirement vector. Its outcome reaches the supremum and weakly dominates every point of
  the open cell. Numerically, in three instances (SKU~1 seed~1; SKU~2 seeds~3 and~6),
  vertices alone and vertices plus probes give identical objective sets (52\,715,
  46\,875 and 67\,989 points). All points that the probe version had marked as
  ``not attained'' are duplicates of attained points.
\end{remark}

By \cref{cor:memo,thm:vertex}, the continuous front is obtained by evaluating $\kappa$ at
the vertices, keeping the largest $B$ per distinct vector, solving one inner problem per
vector, and filtering. In our instances, $\lvert V\rvert$ is several hundred to a few
thousand at $T=48$, and the number of distinct vectors is 7--26\% smaller.

%=====================================================================================
\subsection{Invariance with respect to the ordering cost}\label{app:invariance}

\begin{proposition}[Ordering-cost invariance]\label{prop:c}
  Let \cref{as:lattice} hold and $0<c<p\,r/\Tp$. Then the dominance relation between any
  two candidate solutions (outer point and plan) does not depend on $c$. In particular,
  the Pareto-optimal set of decisions is the same for every $c$ in this interval.
\end{proposition}

\begin{proof}
  $c$ enters only through $\POC=pX(\Tp)+ck$, so it suffices to show that the sign of
  $\Delta\POC=p\,\Delta X+c\,\Delta k$ between two solutions does not depend on $c$. By
  \cref{as:lattice}, $X(\Tp)\in r\mathbb{N}_0$. If $\Delta X\ne0$, then
  $\lvert\Delta X\rvert\ge r$ and $\lvert\Delta k\rvert\le\Tp$, so
  $p\lvert\Delta X\rvert\ge p r>c\Tp\ge c\lvert\Delta k\rvert$, and $\Delta\POC$ has the
  sign of $\Delta X$. If $\Delta X=0$, then $\Delta\POC=c\,\Delta k$ has the sign of
  $\Delta k$ because $c>0$. In both cases the sign, including the case
  $\Delta\POC=0$, is determined independently of $c$.
\end{proof}

\begin{remark}
  For SKU~1 ($p=134$, $r=70$) at $T=48$ ($\Tp=43$) the interval is $0<c<218.1$, which
  contains the values $c=50$ and $c=200$ used in experiment E4. For $c=800$ we also
  observed identical decision sets, which shows that the condition is sufficient but not
  necessary. For SKU~2 ($r=3$) the interval is narrow ($c<33$), and the proposition makes
  no claim at the nominal $c=200$.
\end{remark}

%=====================================================================================
\subsection{Complexity}\label{app:complexity}

Let $S$ be the number of states per layer. With $\hat\beta\approx\mu T$ and spacing $r$
in $\cL$, $S=\Theta(\mu T/r)$. Under \cref{as:lattice}, the partial sum $W$ is a multiple
of $r$ not exceeding $\Tp x_{\max}$, where $x_{\max}$ is the state bound of \cref{lem:ub};
it therefore takes at most $\Tp S$ values, and a node holds labels for at most
$(\Tp+1)\,\Tp S$ distinct pairs $(k,W)$. For each pair only the largest $R$ survives, so
the number $N$ of labels per node satisfies
\begin{equation}
  N=O(\Tp^{2}S).
\end{equation}
A layer has $O(S^2)$ arcs, and each arc extends at most $N$ labels. One inner
problem therefore costs
\begin{equation}\label{eq:complexity}
  O\bigl(\Tp\,S^{2}\,N\log N\bigr)
  =O\bigl(T^{3}S^{3}\log(TS)\bigr)
  =O\bigl(T^{6}(\mu/r)^{3}\log(T\mu/r)\bigr),
\end{equation}
where the logarithm comes from staircase insertion. This is a pseudo-polynomial bound. It
is polynomial in $T$ for a fixed ratio $\mu/r$. The empirical exponents of experiment E3
agree with it: time grows as $T^{6.0\text{--}6.4}$ and as $(\mu/r)^{3.5}$, the stored
labels grow as $T^{4.9\text{--}5.0}$, and time grows as $(\text{labels})^{1.25}$. The
bound is not tight in practice, since \cref{lem:D} removes most ordering arcs and the
staircases remain far below $N$. The number of inner problems is the number of
distinct vectors $\kappa$, which is at most $\lvert V\rvert=O\bigl((\Tp\lvert\Lambda\rvert)^2\bigr)$
for the continuous domain (here $\lvert\Lambda\rvert$ is the number of relevant lattice
values) and at most the grid size in grid mode.

%=====================================================================================
\subsection{Numerical verification of the theoretical results}\label{app:verification}

\Cref{tab:verification} lists, for each result, the experiment that checks it with the
corrected code. Experiment labels refer to \cref{sec:design}.

\begin{table}[h]
  \centering
  \caption{Results of this appendix and their numerical verification.}
  \label{tab:verification}
  \small
  \begin{tabular}{@{}p{3.6cm}p{11.2cm}@{}}
    \toprule
    Result & Verification \\
    \midrule
    \Cref{thm:exact} with \cref{lem:ub,lem:D}
      & E1: DP equals exhaustive enumeration (state bound three times larger, no pruning)
        on 100/100 random instances; at $T=48$ (SKU~2, seed~42) the DP with pruning equals
        the DP without pruning (\cref{rem:c}); all 40 E2 fronts (2 SKUs $\times$ 10 seeds
        $\times$ grid/continuous) are identical before and after the correction. E5a: on 10 + 16 outer points at $T=48$, the
        $\varepsilon$-constraint MILP solved with HiGHS gives identical fronts. \\
    \Cref{thm:vertex}
      & E1: for 20 instances (2 SKUs $\times$ 10 seeds) $\times$ $10^6$ random points of
        $\Omega$ (one quarter on the upper boundary), every point has a vertex with
        $\kappa(v)\le\kappa(z)$ and $B(v)\ge B(z)$; front with and without sector probes
        identical (3 instances); all 20 continuous E2 fronts unchanged after removing the
        probes. \\
    \Cref{prop:c}
      & E4: identical decision sets for $c\in\{50,200\}$ (covered) and $c=800$ (beyond
        the sufficient condition). \\
    Exactness overall
      & E5b: none of the exact points is dominated by any of about $1.2\times10^8$
        NSGA-III evaluations. \\
    Complexity \eqref{eq:complexity}
      & E3: fitted exponents reported above. \\
    \bottomrule
  \end{tabular}
\end{table}


\bibliographystyle{elsarticle-harv}
\bibliography{refs}

\end{document}
